% GENERATED by manuscript/code/build_tables.py -- do not edit by hand.
\begin{table}[t]
\centering
\small
\caption{Per-baseline performance on the temporal holdout ($n=430$) with sequence-cluster bootstrap 95\% CIs and derived calibration states.}
\label{tab:baselines}
\begin{tabular}{lccccc c}
\toprule
Baseline & Brier & 95\% CI & Log-loss & AUROC & Calibration$^a$ & Accuracy \\
\midrule
B0\_no\_skill & 0.232 & [0.212, 0.253] & 0.658 & 0.500 & INFEASIBLE & 0.65 \\
B1\_magnitude\_only\_logit & 0.198 & [0.177, 0.219] & 0.583 & 0.719 & 0.32 / 1.02 & 0.70 \\
B2\_source\_only\_logit & 0.181 & [0.160, 0.202] & 0.536 & 0.786 & 0.64 / 1.29 & 0.73 \\
B4\_random\_forest & 0.187 & [0.168, 0.206] & 0.553 & 0.777 & 0.74 / 1.48 & 0.73 \\
B5\_xgboost & 0.181 & [0.161, 0.202] & 0.541 & 0.783 & 0.42 / 1.07 & 0.74 \\
\bottomrule
\end{tabular}
\par\vspace{2pt}\footnotesize\emph{Note.} $^a$~Intercept / slope. These values cannot be estimated for B0 because it assigns the same probability to every event. B3 includes event year and is therefore omitted from the temporal holdout (Table~\ref{tab:exclusions}). Accuracy uses the probability cutoff of 0.5 fixed on the development set. Lower Brier score and log-loss are better; higher AUROC and accuracy are better. The two lowest-Brier models cannot be distinguished in their direct paired comparison.
\end{table}
